\section*{Statistical Independence}

Statistical independence is the central concept in probability theory that distinguishes probability theory from real analysis. The notion of independent events is closely related to conditional probability. Consider an event $A$ with probability $P(A)$. If we know another event $B$ has already occurred, we can update our knowledge about event $A$ by revising the probability of $A$ to the conditional probability $\frac{P(A \cap B)}{P(B)}$. If events $A$ and $B$ are independent, then the original probability $P(A)$ should equal the conditional probability $\frac{P(A \cap B)}{P(B)}$, meaning that the likelihood of $A$ does not change given the information from event $B$.

This chapter presents a definition of independence of random variables that works for all types of random variables and shows that it is consistent with the definition that is based on probability density function and probability mass function. In the second part of the chapter we discuss the Borel--Cantelli lemmas and their implications.

Throughout this chapter, random variables and events are defined on a probability space $(\Omega, \mathcal{F}, P)$.

\subsection*{5.1 Independence of Two Random Variables}

We will start by defining independence for two events.

\begin{defbox}{5.1}
Two events $A$ and $B$ are said to be \textit{independent} if their joint probability equals the product of their individual probabilities, i.e.,
\[
P(A \cap B) = P(A)P(B).
\]
\end{defbox}

When $P(B)$ is positive, the above definition of independence is equivalent to saying that the conditional probability of $A$ given $B$ is the same as the probability of $A$. However, when $P(B)=0$, both sides of the equations in Definition 5.1 are zero, trivially satisfying the definition of independence.

We can now define independence for random variables in terms of independence of events.

\begin{defbox}{5.2}
Two random variables $X$ and $Y$ defined on the same probability space are said to be \textit{independent} if the inverse images of any two Borel sets $B_1$ and $B_2$ under $X$ and $Y$, respectively, are independent events, i.e., $X^{-1}(B_1)$ and $Y^{-1}(B_2)$ are independent for all Borel sets $B_1$ and $B_2$.

If two random variables $X$ and $Y$ are independent, this implies that the following equation holds for all Borel sets $B_1$ and $B_2$:
\[
P(X \in B_1,\, Y \in B_2) = P(X \in B_1)P(Y \in B_2).
\]
\end{defbox}

Alternatively, we can formulate the independence of two random variables in terms of the $\sigma$-algebras they generate.

\begin{defbox}{5.3}
Given a measurable function $X$ from $(\Omega,\mathcal{F})$ to $(\mathbb{R},\mathcal{B}(\mathbb{R}))$, the $\sigma$-algebra generated by $X$ is defined as
\[
\sigma(X) \triangleq \{X^{-1}(B) : B \in \mathcal{B}(\mathbb{R})\}.
\]

It is easy to check that the collection $\sigma(X)$ of subsets is indeed a sub-$\sigma$-algebra of $\mathcal{F}$. Note that the $\sigma$-algebra generated by $X$ contains all the information about $X$ that is measurable in the sense of Borel sets.
\end{defbox}

\begin{defbox}{5.4}
Two sub-$\sigma$-algebras $\mathcal{F}_1$ and $\mathcal{F}_2$ of a common $\sigma$-algebra $\mathcal{F}$ are said to be \textit{independent} if any set $A_1$ from $\mathcal{F}_1$ and any set $A_2$ from $\mathcal{F}_2$ are independent, i.e.,
\[
P(A_1 \cap A_2) = P(A_1)P(A_2),
\]
for all $A_1 \in \mathcal{F}_1$ and $A_2 \in \mathcal{F}_2$.
\end{defbox}

We can now re-formulate the notion of independence for random variables in terms of the $\sigma$-algebras they generate.

\begin{thmbox}{5.1}
Two random variables $X$ and $Y$ are independent if and only if the $\sigma$-algebras $\sigma(X)$ and $\sigma(Y)$ generated by $X$ and $Y$, respectively, are independent.
\end{thmbox}

If $X$ and $Y$ are independent random variables, it is reasonable to expect that a function of $X$ and a function of $Y$ should also be independent, since they have different sources of randomness. This intuition can be formally justified using the definition of independent random variables.

\begin{thmbox}{5.2}
Suppose $X$ and $Y$ are independent random variables defined on a probability space $(\Omega,\mathcal{F},P)$. If $f$ and $g$ are Borel measurable functions from $\mathbb{R}$ to $\mathbb{R}$, then $f(X)$ and $g(Y)$ are independent random variables.
\end{thmbox}

\textit{Proof} We first note that $f(X(\omega))$ and $g(Y(\omega))$ are both $(\mathcal{F},\mathcal{B}(\mathbb{R}))$-measurable.

To show that $f(X)$ and $g(Y)$ are independent, we need to show that for any two Borel sets $B_1$ and $B_2$ in $\mathbb{R}$, the events
\[
\{\omega : f(X(\omega)) \in B_1\} \text{ and } \{\omega : g(Y(\omega)) \in B_2\}
\]
are independent.

To see this, note that $f^{-1}(B_1)$ and $g^{-1}(B_2)$ are Borel sets in $\mathbb{R}$, because $f$ and $g$ are both Borel measurable. Then, by the independence of $X$ and $Y$, the two events
\[
\{\omega : X(\omega) \in f^{-1}(B_1)\} \text{ and } \{\omega : Y(\omega) \in g^{-1}(B_2)\}
\]
are independent. This shows that $X^{-1}(f^{-1}(B_1))$ and $Y^{-1}(g^{-1}(B_2))$ are independent events, and hence $f(X)$ and $g(Y)$ are independent random variables. \hfill $\square$

Theorem 5.2 is a powerful result that applies to a wide range of functions $f$ and $g$, provided that they are Borel measurable. For example, $f$ and $g$ can be any continuous functions or piece-wise continuous functions. This is because continuous functions and piece-wise continuous functions are Borel measurable.

Moreover, the result in Theorem 5.2 can be generalized to the case where $X$ and $Y$ are independent random vectors. Specifically, if $\mathbf{X}$ and $\mathbf{Y}$ are independent random vectors in $\mathbb{R}^m$ and $\mathbb{R}^n$, respectively, and $\mathbf{f} : \mathbb{R}^m \to \mathbb{R}^p$ and $\mathbf{g} : \mathbb{R}^n \to \mathbb{R}^q$ are Borel measurable functions, then $\mathbf{f}(\mathbf{X})$ and $\mathbf{g}(\mathbf{Y})$ are independent random vectors in $\mathbb{R}^p$ and $\mathbb{R}^q$, respectively. The proof is similar to the proof of Theorem 5.2.

The theory we develop in this section works for random variables of any types, including discrete, continuous, singular, or mixed-type random variables. This framework is sufficiently general to address the independence of a discrete random variable and a continuous random variable.

\textbf{Example 5.1.1 (Independent Random Variables of Two Different Types)} \\
We randomly select a point uniformly from the unit square $[0,1] \times [0,1]$ and let the random point be $(X,Y)$, where $X$ and $Y$ are independent uniform random variables. We define random variable $U$ as $\lfloor 10X \rfloor$ and random variable $V$ as $-\ln(Y)$. Random variable $U$ is discrete, while random variable $V$ is continuous. Although the joint distribution function of $U$ and $V$ is undefined, we can still conclude that $U$ and $V$ are independent random variables.
